\section{Introduction}
\label{sec:introduction}

Automated biomedical image segmentation is fundamental to computer-assisted diagnosis, radiotherapy planning, and surgical navigation~\cite{ronneberger2015unet,zhou2018unetpp,chen2021transunet}. Clinical modalities span a broad spectrum---from 3D volumetric CT and MRI to 2D endoscopic video and high-resolution dermoscopy. Accurate boundary delineation across these domains requires architectures capable of both fine-grained edge localization and expansive global context modeling.

While convolutional neural networks (CNNs)~\cite{ronneberger2015unet,zhou2018unetpp,oktay2018attention} form the conventional baseline, their localized receptive fields limit long-range context aggregation. Hybrid CNN-Transformer frameworks like TransUNet~\cite{chen2021transunet} address this via self-attention, and dual-attention networks~\cite{fu2019dual} concurrently capture spatial affinities via Position Attention Modules (PAM) and cross-channel correlations via Channel Attention Modules (CAM). Recently, DA-TransUNet~\cite{sun2024datransunet} integrated dual-attention blocks across multi-scale skip connections and the bottleneck, demonstrating state-of-the-art accuracy on multi-organ CT and colon polyp benchmarks.

Despite their representational strength, standard dual-attention architectures are computationally heavy. Full-map PAM scales quadratically with the number of spatial tokens and CAM with the number of channels, and DA-TransUNet couples these modules with a ViT-B/16 encoder and a wide decoder, which brings it above 100M parameters and makes it costly to train and deploy on resource-constrained hardware. The quadratic cost of attention has prompted efficient approximations---windowed attention~\cite{liu2021swin}, low-rank factorization~\cite{wang2020linformer}, and grouped channel interaction~\cite{rahman2024emcad}---but in DA-TransUNet the dual-attention blocks account for less than $3\%$ of the parameters and FLOPs; most of the cost lies in the ViT-B/16 encoder and the wide decoder around them.

In our preliminary workshop study~\cite{zhong2026approxda}, we introduced the ApproxDA block, which keeps the DANet head of DA-TransUNet and replaces its attention operators with three independently controllable approximations: windowed position attention (window size $M$), low-rank projection (rank $r$), and grouped channel attention (groups $G$). In that study, exact attention showed no consistent advantage over its low-rank approximation, but the model kept the R50-ViT-B/16 backbone, so the savings of the approximated attention barely reached the whole network. In this work, we carry the approximated dual attention into a lightweight architecture. We present \textbf{ApproxDA-TransUNet}, a system-level redesign of DA-TransUNet built from four components:
(i) a \emph{ResNet-50 stem truncated after its second stage} followed by a compact Vision Transformer (\emph{LightViT-T}~\cite{huang2022lightvit}), which makes the encoder about ten times smaller;
(ii) \emph{approximated dual-attention blocks}, in which position attention is restricted to local windows and projected to a low rank, channel attention is computed in groups, and the two branches are blended with a fixed, equal gate;
(iii) these blocks on all three skip connections and on the bottleneck after the transformer, whereas DA-TransUNet applies dual attention to two of the skips and to the patch embeddings before the ViT; and
(iv) a \emph{slim decoder} with additive skip fusion and narrow channels at every skip stage.
Together, these changes reduce the parameters to ${\sim}10\%$, the FLOPs to ${\sim}23\%$, and the GPU memory for training and inference to ${\sim}42\%$ of those of DA-TransUNet.

In single-axis ablations on Synapse, removing the dual-attention blocks costs nearly three Dice points, whereas their approximation costs no accuracy: the exact global dual attention of DA-TransUNet gives the same Dice with more computation and twice the training time.

We evaluate ApproxDA-TransUNet with three seeds on five benchmarks spanning CT, cardiac MRI, colonoscopy, dermoscopy, and fetal brain MRI: Synapse, ACDC, CVC-ClinicDB, ISIC 2018, and FeTA 2.4. Its mean Dice is higher than that of DA-TransUNet on all five datasets, by $0.37$ points on average; the gain on any single dataset is within seed variation. On FeTA 2.4, it is about two points above the reproduced U-Net family and Swin-Unet and on par with TransUNet, at a tenth of TransUNet's size.

The primary contributions of this paper are:
\begin{itemize}
    \item \textbf{Lightweight dual-attention architecture}: ApproxDA-TransUNet retains the dual attention of DA-TransUNet with ${\sim}10\%$ of its parameters, ${\sim}23\%$ of its FLOPs, and ${\sim}42\%$ of its GPU memory, at comparable segmentation accuracy, making dual-attention segmentation practical on modest clinical hardware.
    \item \textbf{Dual-attention approximation}: windowed low-rank position attention ($M$, $r$) and grouped channel attention ($G$) as three independently parameterized axes, with single-axis ablations on Synapse showing that the blocks are needed and that the approximation matches exact dual attention at lower cost.
    \item \textbf{Five-dataset, three-seed validation}: paired comparisons with DA-TransUNet on benchmarks spanning CT, MRI, colonoscopy, dermoscopy, and fetal MRI, with literature and reproduced baselines and a deployment profile.
\end{itemize}

The rest of this paper is organized as follows: Section~\ref{sec:related_work} reviews related work; Section~\ref{sec:method} presents the methodology; Section~\ref{sec:experiments} reports experimental results; Section~\ref{sec:ablation} provides ablation studies; Section~\ref{sec:discussion} discusses augmentation, deployment cost, and limitations; Section~\ref{sec:conclusion} concludes the paper.
